Durante el período 2018–2025 las empresas analizadas mostraron diferencias significativas tanto en su rendimiento acumulado como en su nivel de riesgo. NVIDIA obtuvo el mayor rendimiento acumulado del conjunto, mientras que Tesla y BioNTech también registraron crecimientos sobresalientes, aunque acompañados de una mayor volatilidad.
La evolución de los precios de cierre evidenció que un alto rendimiento no necesariamente implica un comportamiento altamente inestable. Mientras BNTX experimentó fuertes incrementos seguidos de importantes correcciones, NVIDIA mantuvo una trayectoria de crecimiento más sostenida y consistente durante el período analizado.
Los análisis de volatilidad, rendimientos diarios y distribuciones de densidad mostraron que BNTX, TSLA y PDD fueron las empresas con mayor dispersión en sus retornos, mientras que compañías como SNPS, KLAC, CDNS, FICO y ARES registraron comportamientos considerablemente más estables.
El volumen de negociación permitió identificar a Tesla, Apple y AMD como las empresas con mayor liquidez del mercado. Sin embargo, no se encontró una relación lineal clara entre el volumen negociado y el precio de cierre, lo que indica que ambas variables responden a factores distintos del comportamiento bursátil.
El análisis mensual evidenció ciertos patrones estacionales compartidos entre las empresas, observándose meses con rendimientos promedio predominantemente positivos y otros con resultados negativos. No obstante, la magnitud de dichos cambios varía considerablemente entre compañías, reforzando la importancia de analizar tanto el rendimiento esperado como el nivel de riesgo asociado.